\documentclass[12pt, a4paper]{article}

% --- Packages ---
\usepackage[utf8]{inputenc}
\usepackage[T1]{fontenc}
\usepackage{geometry}
\geometry{left=2.5cm, right=2.5cm, top=2.5cm, bottom=2.5cm}
\usepackage{times}
\usepackage[numbers,sort&compress]{natbib}
\usepackage{hyperref}
\usepackage{graphicx}
\usepackage{titlesec}
\usepackage{enumitem}
\usepackage{setspace}
\usepackage{amsmath}
\usepackage{amssymb}
\usepackage{booktabs}
\usepackage{array}
\usepackage{multirow}
\usepackage{xcolor}

\hypersetup{
    colorlinks=true,
    linkcolor=blue,
    filecolor=magenta,
    urlcolor=cyan,
    citecolor=blue,
    pdfencoding=auto
}

\setstretch{1.15}
\setlength{\parskip}{0.35em}
\titlespacing{\section}{0pt}{12pt}{5pt}
\titlespacing{\subsection}{0pt}{8pt}{3pt}
\titlespacing{\subsubsection}{0pt}{6pt}{2pt}

% =========================================================
\title{\textbf{Multi-Agent Reinforcement Learning for Autonomous Driving:
A Condensed Survey of Algorithms, Coordination, and Cognitive Paradigms}}
\author{\large Literature Review}
\date{\today}

\begin{document}

\maketitle

% =========================================================
\begin{abstract}
Autonomous driving is fundamentally a multi-agent decision problem: every plan
for the ego vehicle depends on how surrounding agents react, and those reactions
in turn depend on the ego vehicle's own future behavior. This survey reviews the
field through three linked paradigms. The first covers single-agent learning
foundations, including end-to-end imitation learning, covariate-shift mitigation,
and early reinforcement learning for driving. The second focuses on the core
multi-agent turn: Centralized Training with Decentralized Execution (CTDE),
value decomposition, actor-critic coordination, communication protocols,
interaction-aware prediction, and cooperative driving applications. The third
surveys recent scalable and cognitive architectures, especially Transformer-based
multi-agent policies and large language model (LLM) driven driving agents that
inject commonsense reasoning, memory, and neuro-symbolic control into the stack.

The survey emphasizes field-shaping methods and benchmark-defining systems while
organizing a broader set of representative studies into compact thematic
comparisons. This keeps the paper survey-like rather than encyclopedic while
preserving coverage of more than 100 references. Across paradigms, three
tensions recur: expressiveness versus
tractability in credit assignment, explicit communication versus implicit
coordination, and closed-world optimization versus robustness under long-tail
shift. We conclude that the most promising path forward is neither purely
end-to-end nor purely symbolic, but hybrid: interaction-aware MARL and world
models at the policy level, coupled with structured safety constraints and
higher-level reasoning modules for rare, ambiguous traffic scenarios.
\end{abstract}

\noindent\textbf{Keywords:} Autonomous Driving, Multi-Agent Reinforcement Learning,
Imitation Learning, Centralized Training Decentralized Execution, Graph Neural
Networks, Cooperative Perception, Large Language Models, Safety, Simulation.

% ==========================================================
\section{Introduction}
\label{sec:intro}

Autonomous driving (AD) requires robust perception, intent prediction,
interactive planning, and low-latency control in scenes populated by multiple
vehicles, pedestrians, and infrastructure entities. The resulting problem is not
well captured by a purely single-agent view: surrounding actors are strategic,
partially observed, and behaviorally coupled to the ego vehicle. For that
reason, the historical development of learning-based AD is best understood as a
progression from \emph{single-agent policy learning}, to \emph{explicit
multi-agent coordination}, and then to \emph{cognitive and scalable
architectures} that try to generalize beyond closed training distributions.

The survey is organized around that progression. Paradigm~I established the
foundations with end-to-end behavior cloning and early driving RL, but remained
fragile under covariate shift and interaction misspecification. Paradigm~II
introduced the main MARL toolkit: CTDE, value decomposition, actor-critic
methods, learned communication, graph-based interaction modeling, and dedicated
multi-agent simulators. Paradigm~III extends the field toward Transformer-based
sequence models and LLM-augmented agents that can reason over rare or ambiguous
scenarios, albeit with unresolved latency, grounding, and safety issues.

\subsection{Scope and Selection}
\label{sec:scope}

This condensed survey retains the full thematic breadth of the longer draft but
compresses exposition by treating canonical papers in detail and covering closely
related follow-up work comparatively. The literature base spans more than 100
references from foundational game-theoretic MARL to contemporary LLM-based AD,
with emphasis on methodological turning points, benchmark-defining systems, and
studies that exposed limitations later addressed by the next paradigm.

We retain four survey goals. First, we place AD consistently in the POMDP/POSG
framework. Second, we compare algorithm families rather than listing papers one
by one. Third, we connect representation learning, policy optimization,
communication, and evaluation infrastructure under a shared multi-agent lens.
Fourth, we include the emerging cognitive-agent literature rather than stopping
at classical MARL. The resulting structure follows a mainstream survey pattern:
problem formulation, taxonomy, paradigm-wise review, cross-cutting synthesis,
and open problems.

\subsection{Representative Literature Map}
\label{sec:litmap}

To keep the narrative concise without sacrificing breadth, we summarize the
reviewed literature by theme below.

\begin{itemize}[leftmargin=*,itemsep=0.25em,topsep=0.3em]
  \item \textbf{Foundational MARL theory and general surveys:}
  stochastic and Markov game formulations, equilibrium analysis, and broad MARL
  overviews appear in \cite{shapley1953stochastic,littman1994markov,bernstein2002complexity,hu2003nash,busoniu2008comprehensive,lanctot2017unified,zhang2021multi,hernandez2019survey,tan1993multi,kraemer2016multi,foerster2018lola,vinyals2019alphastar,gupta2017cooperative,zhang2019global,de2020independent}.

  \item \textbf{Single-agent AD learning foundations:}
  end-to-end imitation and single-agent RL are represented by
  \cite{bojarski2016end,codevilla2018end,bansal2018chauffeur,ross2011reduction,sallab2017deep,mnih2015human,mnih2016asynchronous,schulman2017proximal,kahn2017uncertainty,chen2019model,kiran2021deep,hubmann2018automated}.

  \item \textbf{Interaction modeling, prediction, and planning:}
  social pooling, graph-based prediction, and planning-aware forecasting include
  \cite{alahi2016social,gupta2018social,lee2017desire,deo2018convolutional,gao2020vectornet,salzmann2020trajectron,liang2020lanegcn,zhao2020tnt,chai2019multipath,song2020pip,sadigh2016planning,fisac2019hierarchical,schwarting2019social,schwarting2021deep,fridovich2020efficient,deo2020trajectory}.

  \item \textbf{Core MARL algorithm families:}
  value decomposition, actor-critic CTDE, role-based, and hierarchical methods
  include \cite{sunehag2018value,rashid2018qmix,son2019qtran,mahajan2019maven,lowe2017multi,foerster2018counterfactual,yu2021surprising,kuba2021hatrpo,peng2021facmac,wang2021rode,wang2020roma,vezhnevets2017feudal,chu2020multi,wang2020shapley,su2021vdac}.

  \item \textbf{Communication, coordination, and cooperative perception:}
  differentiable communication and V2X-oriented methods include
  \cite{foerster2016learning,sukhbaatar2016learning,das2019tarmac,jiang2018learning,singh2019learning,kim2019learning,kim2019message,wang2021vbc,wang2020bandwidth,wang2020v2vnet,jiang2019graph,velickovic2018graph}.

  \item \textbf{Applications and mixed-autonomy traffic control:}
  representative work on traffic signal control, ramp merging, and connected
  traffic systems includes
  \cite{wei2019colight,wei2019presslight,zheng2019learning,wang2018urban,wu2017flow,stern2018dissipation,talebpour2016influence,isele2018navigating,shalev2016safe,palanisamy2020multi}.

  \item \textbf{Simulation platforms and datasets:}
  the main evaluation ecosystem includes
  \cite{dosovitskiy2017carlaopenurbandriving,krajzewicz2012recent,zhou2020smarts,caesar2020nuscenes,caesar2021nuplan,chang2019argoverse,ettinger2021large,waymo2021open,samvelyan2019starcraft,papoudakis2021benchmarking}.

  \item \textbf{Transformers, LLM agents, safety, and transfer:}
  the recent scalable and cognitive literature includes
  \cite{vaswani2017attention,wen2022multi,chen2021decision,janner2021offline,parisotto2020stabilizing,brown2020language,fu2023drive,wen2023dilu,mao2023gpt,sha2023languagempc,wang2023voyager,achiam2017constrained,alshiekh2018safe,prashanth2016variance,tobin2017domain,finn2017model,he2016opponent,zhang2018fully,li2020deep}.
\end{itemize}

\subsection{Three-Layer Taxonomy}
\label{sec:taxonomy}

The rest of the survey uses a three-layer taxonomy. \textbf{Layer~1} covers
single-agent learning foundations: behavior cloning, imitation learning, and
early driving RL. \textbf{Layer~2} covers explicit multi-agent coordination:
MARL algorithms, communication mechanisms, and interaction-aware scene models.
\textbf{Layer~3} covers scalable and cognitive architectures: Transformer-based
policies, LLM planners, and hybrid neuro-symbolic control. The layers are not
mutually exclusive; modern systems often combine a Layer~2 scene encoder,
Layer~2 MARL policy, and Layer~3 reasoning module. Their value lies in
structuring the field's main shifts in emphasis.

\begin{table}[ht]
\centering
\caption{Three paradigm shifts in learning-based autonomous driving.
\label{tab:paradigm_shift}}
\renewcommand{\arraystretch}{1.2}
\small
\begin{tabular}{p{2.2cm}p{3.0cm}p{3.2cm}p{3.2cm}p{2.6cm}}
\toprule
\textbf{Paradigm} & \textbf{Problem that became dominant} & \textbf{Main innovation} & \textbf{Landmark studies} & \textbf{Main residual issue} \\
\midrule
I. Single-agent foundations & Covariate shift, modular error propagation, weak interaction modeling & End-to-end imitation, conditional imitation, early policy optimization & \cite{bojarski2016end,bansal2018chauffeur,codevilla2018end,schulman2017proximal} & Other agents treated as environment \\
\addlinespace
II. MARL and coordination & Non-stationarity, credit assignment, communication, cooperative control & CTDE, value decomposition, actor-critic MARL, graph-based and communication-aware models & \cite{lowe2017multi,rashid2018qmix,yu2021surprising,das2019tarmac,gao2020vectornet,wang2020v2vnet} & Scale, long-tail robustness, interpretability \\
\addlinespace
III. Scalable and cognitive agents & Sparse long-tail coverage, opaque decisions, weak semantic reasoning & Transformer policies, sequence modeling, LLM reasoning, neuro-symbolic hybrids & \cite{wen2022multi,fu2023drive,wen2023dilu,mao2023gpt,sha2023languagempc} & Latency, grounding, verification, sim-to-real gap \\
\bottomrule
\end{tabular}
\end{table}

% ==========================================================
\section{Problem Formulation and Background}
\label{sec:background}

\subsection{From POMDP to POSG}

For a single vehicle, the driving problem is typically formalized as a partially
observable Markov decision process (POMDP), with hidden world state
$s_t \in \mathcal{S}$, action $a_t \in \mathcal{A}$, observation
$o_t \in \mathcal{O}$, transition model $T$, reward $R$, and discount factor
$\gamma$. The objective is
\begin{equation}
\pi^* = \arg\max_{\pi} \; \mathbb{E}_{\tau \sim \pi}
\left[\sum_{t=0}^{T} \gamma^t R(s_t,a_t)\right].
\label{eq:pomdp}
\end{equation}
This formalism is useful for end-to-end imitation learning and single-agent RL,
but it is incomplete for interactive traffic, because the transition dynamics are
jointly determined by many agents.

A more faithful formalization is the partially observable stochastic game (POSG),
which introduces an agent set $\mathcal{N}$, per-agent observations and actions,
and a joint transition model
\begin{equation}
T(s_{t+1} \mid s_t, a_t^1, \dots, a_t^N).
\label{eq:posg}
\end{equation}
This shift matters because the core pathology of MARL is \emph{non-stationarity}:
from the perspective of any one learner, the environment changes when the other
agents' policies change \cite{littman1994markov,bernstein2002complexity,lowe2017multi,zhang2021multi}. In traffic, rewards are usually mixed rather than purely cooperative or purely adversarial: agents share collision-avoidance incentives but compete for space, right-of-way, or short-term efficiency.

\subsection{Why AD Stress-Tests MARL}

Autonomous driving is an unusually hard MARL domain for four reasons. First,
observations are partial and noisy, and safety-relevant latent variables such as
human intent are never directly observed. Second, action spaces are often
continuous, which weakens the applicability of purely discrete MARL methods.
Third, rare failures dominate practical importance: a method that works on common
scenes but fails under merges, occlusions, or ambiguous negotiation is not
deployment-ready. Fourth, evaluation is expensive and fragmented across
simulators, logged datasets, and closed-loop tests \cite{caesar2021nuplan,zhou2020smarts,dosovitskiy2017carlaopenurbandriving,kiran2021deep}.

\subsection{Software-Stack View}

The learning problem can sit at several locations in the AD stack. Some systems
learn a direct observation-to-control policy \cite{bojarski2016end}. Others keep
a modular pipeline and learn only prediction or planning components
\cite{gao2020vectornet,liang2020lanegcn,song2020pip}. Multi-agent work is also
split between \emph{vehicle-level decision making} (merging, negotiation,
cooperative driving), \emph{infrastructure coordination} (traffic signal
control), and \emph{mixed-autonomy traffic shaping}
\cite{wei2019colight,wei2019presslight,stern2018dissipation,wu2017flow}.
Keeping these deployment locations separate is crucial when comparing methods,
because what counts as ``state'', ``action'', and even ``agent'' varies sharply
across them.

% ==========================================================
\section{Paradigm I: Single-Agent Learning Foundations}
\label{sec:paradigm1}

\subsection{Behavior Cloning and Its Limits}

The early modern wave of learning-based AD was driven by end-to-end imitation.
DAVE-2 demonstrated that a convolutional network could map front-camera images
directly to steering commands \cite{bojarski2016end}. Conditional imitation
learning later improved route compliance by conditioning on high-level commands
such as turn-left or go-straight \cite{codevilla2018end}. These systems were
important because they showed that handcrafted intermediate representations were
not strictly necessary for limited driving tasks.

Their central weakness was covariate shift. Supervised training on expert data
concentrates probability mass on nominal states; once the learned policy drifts,
it enters states absent from the demonstrations and has no learned recovery
behavior. The classic reduction from imitation to online learning in DAgger
formalized this issue \cite{ross2011reduction}. In AD, ChauffeurNet was the most
influential practical response: it used perturbation-based synthesis to expose
the model to off-trajectory and near-failure states during training, showing that
scale alone does not solve robustness \cite{bansal2018chauffeur}. This paper is
arguably the clearest end point of purely single-agent imitation for driving:
stronger than naive behavior cloning, but still fundamentally limited by how
other agents are modeled.

\subsection{Single-Agent RL for Driving}

Reinforcement learning offered a way to optimize long-horizon safety and
efficiency rather than imitating expert trajectories. General deep RL advances
such as DQN, asynchronous actor-critic methods, and PPO supplied the algorithmic
basis \cite{mnih2015human,mnih2016asynchronous,schulman2017proximal}. In AD,
representative single-agent formulations include uncertainty-aware collision
avoidance and model-based policy learning under limited data
\cite{kahn2017uncertainty,chen2019model}, along with broader surveys of driving
RL \cite{sallab2017deep,kiran2021deep}.

However, once traffic becomes interactive, the single-agent abstraction begins to
fail for structural rather than merely data-efficiency reasons. The behavior of
surrounding vehicles is not an exogenous disturbance but a strategic response to
the ego plan. Treating other agents as part of the environment obscures joint
credit assignment and induces unstable learning when multiple adaptive actors are
present \cite{zhang2021multi,busoniu2008comprehensive}.

\subsection{Early Interaction Modeling Before Full MARL}

Before explicit MARL became dominant, the field moved through an intermediate
phase that improved prediction without fully solving coordination. Social LSTM
showed that neighboring agents should be modeled jointly rather than
independently \cite{alahi2016social}. Later work replaced rasterized pooling with
richer scene representations: DESIRE and convolutional social pooling improved
multi-modal trajectory prediction \cite{lee2017desire,deo2018convolutional},
while graph- and vector-based models such as Trajectron++, VectorNet, LaneGCN,
TNT, and MultiPath captured lane topology and long-range interaction more
explicitly \cite{gupta2018social,salzmann2020trajectron,gao2020vectornet,liang2020lanegcn,zhao2020tnt,chai2019multipath}.

The strongest conceptual bridge to Paradigm~II was planning-aware interaction
modeling. PiP argued that prediction should be conditioned on candidate ego
plans, because other agents react to what the ego vehicle intends to do
\cite{song2020pip}. Related work on socially compatible motion planning and
hierarchical game-theoretic planning made the same point from a control
perspective \cite{sadigh2016planning,fisac2019hierarchical,schwarting2019social,schwarting2021deep,fridovich2020efficient}. These methods significantly improved interaction reasoning, but they still left coordination mostly implicit rather than jointly optimized.

\subsection{What Paradigm~I Could Not Resolve}

Paradigm~I established the key ingredients of learning-based AD but left three
open problems. First, other agents remained only partially modeled, so the
non-stationarity of genuine multi-agent learning was not addressed. Second,
credit assignment across interacting vehicles was absent. Third, many methods
optimized open-loop prediction quality rather than closed-loop joint behavior.
These limitations directly motivated the explicit MARL turn.

% ==========================================================
\section{Paradigm II: MARL Algorithms and Coordination}
\label{sec:paradigm2}

\subsection{CTDE as the Central Organizing Principle}

The most important conceptual advance of Paradigm~II was centralized training
with decentralized execution (CTDE). The idea is simple but powerful: during
training, a critic or mixing function can access global information; during
execution, each agent acts from local observations only \cite{kraemer2016multi,lowe2017multi}. CTDE solves neither credit assignment nor scalability by itself,
but it gives MARL a workable interface between learning and deployment.

Two algorithm families became especially influential. \emph{Value decomposition}
methods factorize a joint value function into per-agent terms. VDN enforced an
additive decomposition, while QMIX introduced a monotonic mixing network that
preserves decentralized argmax consistency and became the canonical baseline for
cooperative MARL \cite{sunehag2018value,rashid2018qmix}. QTRAN and MAVEN tried
to relax the expressiveness bottlenecks of QMIX, although at greater complexity
and with less clear practical dominance \cite{son2019qtran,mahajan2019maven}.

\emph{Actor-critic CTDE} methods addressed continuous control more naturally.
MADDPG showed that centralized critics could stabilize multi-agent policy
gradients in mixed settings \cite{lowe2017multi}. COMA introduced a
counterfactual advantage for exact cooperative credit assignment in discrete
action spaces \cite{foerster2018counterfactual}. MAPPO later demonstrated that a
comparatively simple centralized-value PPO baseline is surprisingly competitive
across benchmarks, making it a highly relevant practical reference point
\cite{yu2021surprising}. HATRPO and FACMAC are important specialized extensions:
HATRPO adds stronger optimization guarantees, while FACMAC improves sample
efficiency for high-cost simulators \cite{kuba2021hatrpo,peng2021facmac}.

Other variants are best understood comparatively rather than as independent
turning points. Independent baselines remain useful stress tests
\cite{tan1993multi,de2020independent}. Role-based methods encode structured
heterogeneity \cite{wang2021rode,wang2020roma}; hierarchical methods improve
temporal abstraction \cite{vezhnevets2017feudal,chu2020multi}; and recent
credit-assignment work refines cooperative value factorization
\cite{wang2020shapley,su2021vdac}. Together, these works broaden the design
space but do not displace QMIX, MADDPG, and MAPPO as the main reference points.

\subsection{Communication and Coordination Protocols}

When coordination cannot emerge reliably from reward structure alone,
communication becomes explicit. Differentiable communication began with recurrent
message-passing approaches and mean-pooling baselines \cite{foerster2016learning,sukhbaatar2016learning}. Later methods learned \emph{who} should communicate
with \emph{whom}. TarMAC introduced targeted communication with attention, while
other work learned communication schedules, gating, or compressed messages
\cite{das2019tarmac,jiang2018learning,singh2019learning,kim2019learning,kim2019message}.

In driving, the communication problem is more constrained than in generic MARL.
Real systems operate under bandwidth, latency, packet loss, and sensing-range
limits. As a result, V2X-oriented work matters more than generic message passing
alone. V2VNet is representative because it moves communication into the feature
space of cooperative perception, which is closer to realistic connected driving
than symbol-level messages between abstract agents \cite{wang2020v2vnet}. Related
work on bandwidth control and viewpoint-aware broadcasting makes the same point:
communication quality is itself a resource-allocation problem
\cite{wang2021vbc,wang2020bandwidth}. Graph-attentional communication also
bridges generic MARL and structured traffic interaction \cite{jiang2019graph,velickovic2018graph}.

The broader lesson is that communication is most useful when it conveys
information unavailable to local sensing, not when it simply duplicates what a
well-designed centralized critic already exploited during training.

\subsection{Interaction-Aware Representation Learning}

A second line of Paradigm~II made interaction structure explicit in the state
representation. Graph neural networks and vectorized map encodings became the
default choice because traffic scenes are relational rather than grid-like.
VectorNet and LaneGCN were especially influential in replacing heavy raster
representations with structured polyline graphs \cite{gao2020vectornet,liang2020lanegcn}. Trajectron++, TNT, and MultiPath pushed multi-modal prediction and
scene-conditioned forecasting further \cite{salzmann2020trajectron,zhao2020tnt,chai2019multipath}. These models are not MARL algorithms by themselves, but in
practice they became the dominant front-end for MARL-style decision systems.

The key conceptual advance was not merely better prediction accuracy; it was the
recognition that prediction and planning should be tightly coupled. PiP,
socially compatible planning, and hierarchical interaction reasoning all showed
that forecasting other agents independently of the ego plan produces brittle
closed-loop behavior \cite{song2020pip,sadigh2016planning,fisac2019hierarchical,schwarting2021deep}. This insight narrowed the gap between trajectory prediction and joint policy learning.

\subsection{Applications: Traffic Control, Mixed Autonomy, and Cooperative Driving}

Paradigm~II matured because it was tested on concrete application classes rather
than only toy benchmarks. In traffic signal control, graph-based MARL systems
such as CoLight and PressLight showed that coordination among intersections can
be framed as a scalable cooperative control problem \cite{wei2019colight,wei2019presslight,zheng2019learning,wang2018urban}. In mixed-autonomy traffic,
Flow and related studies demonstrated that even a small fraction of learned or
controlled vehicles can improve macro-scale traffic efficiency and dampen stop-
and-go waves \cite{wu2017flow,stern2018dissipation,talebpour2016influence}.
Vehicle-level navigation and highway interaction studies further extended MARL to
merging, lane changes, and risk-sensitive negotiation
\cite{isele2018navigating,hubmann2018automated,shalev2016safe,palanisamy2020multi}.

These applications revealed an important asymmetry. Traffic signal control often
fits cooperative MARL cleanly because agent interfaces are explicit and the
environment is relatively structured. Vehicle-level AD is harder: sensing is
partial, action spaces are continuous, and the cost of rare failure is much
higher. Results therefore transfer imperfectly across application classes even
when the algorithmic vocabulary looks similar.

\subsection{Evaluation Ecosystem}

Progress in Paradigm~II depended heavily on simulators and benchmarks. CARLA
provided high-fidelity urban simulation with perception stacks
\cite{dosovitskiy2017carlaopenurbandriving}; SUMO enabled cheap large-scale traffic
experiments \cite{krajzewicz2012recent}; SMARTS and MACAD-Gym made multi-agent AD
training more systematic \cite{zhou2020smarts,palanisamy2020multi}; and SMAC,
though not an AD benchmark, remained a standard cooperative MARL stress test
\cite{samvelyan2019starcraft,papoudakis2021benchmarking}. On the data side,
nuScenes, Argoverse, Waymo Open, and WOMD expanded perception and motion
forecasting coverage, while nuPlan pushed closed-loop planning evaluation closer
to deployment needs \cite{caesar2020nuscenes,chang2019argoverse,waymo2021open,ettinger2021large,caesar2021nuplan}.

\begin{table}[ht]
\centering
\caption{Compact comparison of widely used simulation and dataset ecosystems.
\label{tab:platforms}}
\renewcommand{\arraystretch}{1.15}
\small
\begin{tabular}{p{2.3cm}p{1.3cm}p{1.8cm}p{1.6cm}p{5.6cm}}
\toprule
\textbf{Platform} & \textbf{Type} & \textbf{Closed-loop} & \textbf{Native MARL} & \textbf{Typical role in the literature} \\
\midrule
CARLA \cite{dosovitskiy2017carlaopenurbandriving} & Simulator & Yes & Partial & High-fidelity perception-to-control and urban interaction studies \\
SUMO \cite{krajzewicz2012recent} & Simulator & Yes & Yes & Large-scale traffic control and mixed-autonomy experiments \\
SMARTS \cite{zhou2020smarts} & Simulator & Yes & Yes & Standardized multi-agent AD scenarios and benchmarking \\
SMAC \cite{samvelyan2019starcraft} & Simulator & Yes & Yes & Non-AD cooperative MARL baseline development \\
nuScenes \cite{caesar2020nuscenes} & Dataset & No & No & Perception and motion-prediction evaluation \\
Argoverse \cite{chang2019argoverse} & Dataset & No & No & Trajectory forecasting and scene understanding \\
WOMD \cite{ettinger2021large} & Dataset & Partial & No & Large-scale motion forecasting and logged interaction data \\
nuPlan \cite{caesar2021nuplan} & Benchmark & Yes & No & Closed-loop planning and planner comparison \\
Waymo Open \cite{waymo2021open} & Dataset & No & No & Large-scale multi-sensor supervision and evaluation support \\
\bottomrule
\end{tabular}
\end{table}

\subsection{Why Paradigm~II Still Plateaued}

Paradigm~II solved the language of coordination but not the long-tail problem.
Three limitations became increasingly visible. First, algorithmic improvements
were often benchmark-sensitive: methods that excelled in cooperative game suites
did not necessarily dominate in closed-loop driving. Second, many communication
and CTDE setups assumed cleaner observability or tighter synchronization than
real V2X systems can guarantee. Third, even strong interaction-aware MARL agents
remained opaque and brittle under rare or semantically unusual situations. These
limitations opened the door to Paradigm~III.

% ==========================================================
\section{Paradigm III: Scalable and Cognitive Architectures}
\label{sec:paradigm3}

\subsection{Transformer-Based Policies and Sequence Models}

Transformers entered MARL because attention scales more gracefully to variable
agent sets and long contexts than recurrent encoders. The most representative
multi-agent example in this survey is MAT, which models joint action generation
autoregressively and showed that sequence modeling can compete strongly with
classical value-decomposition and actor-critic approaches \cite{wen2022multi}. The
broader trend is supported by general sequence-modeling work such as Transformer
attention, Decision Transformer, and trajectory-sequence modeling
\cite{vaswani2017attention,chen2021decision,janner2021offline,parisotto2020stabilizing}.

For AD, the appeal of Transformer-based policies is twofold. First, traffic
scenes naturally consist of sets or sequences of agents, lanes, and interaction
histories. Second, attention provides a more interpretable handle on which actors
matter at a given moment. Yet sequence models do not automatically solve safety
or grounding; they mostly provide a stronger architecture for relational context
aggregation. Their practical value is highest when paired with structured scene
encoders and disciplined training objectives.

\subsection{LLMs as Cognitive Driving Agents}

The most visible shift of Paradigm~III is the use of LLMs as high-level driving
reasoners. Drive-like prompting frameworks showed that large language models can
produce stepwise scene interpretation and decision rationales in novel settings
\cite{fu2023drive,brown2020language}. DiLu added reflection and memory for
lifelong adaptation \cite{wen2023dilu}. GPT-Driver treated motion planning as a
language-generation problem \cite{mao2023gpt}. LanguageMPC proposed a more
practical hybrid pattern: use language reasoning for strategic decomposition or
cost shaping, but retain model predictive control for low-level execution
\cite{sha2023languagempc}. Work such as Voyager is also relevant because it shows
how LLM agents can accumulate reusable skills through memory and self-improvement
loops \cite{wang2023voyager}.

At present, these systems are best interpreted as \emph{cognitive overlays}
rather than replacements for continuous control. Their strengths are semantic
reasoning, explanation, and adaptation to rare scenarios; their weaknesses are
latency, susceptibility to hallucination, dependence on prompt design, and lack
of formal guarantees. For AD, the most credible near-term use is therefore a
hybrid stack in which the LLM operates at the tactical or deliberative level,
while classical planners or constrained policies execute the final control.

\subsection{Safety, Transfer, and Scalability as Cross-Cutting Constraints}

Paradigm~III also revived interest in constraints that classical benchmark-driven
MARL often underemphasized. Safe RL and constrained optimization matter because
rare collisions dominate deployment risk \cite{achiam2017constrained,alshiekh2018safe,prashanth2016variance,shalev2016safe}. Sim-to-real transfer remains difficult
because simulators underrepresent long-tail behavior and sensing artifacts
\cite{tobin2017domain,finn2017model}. Opponent modeling and explicit adaptation
remain relevant for mixed traffic with humans or heterogeneous autonomous agents
\cite{he2016opponent,zhang2018fully}. Finally, scaling to large populations or
variable agent counts requires architectures that avoid quadratic coordination
costs, motivating attention-based and deep-set style abstractions
\cite{li2020deep,wen2022multi}.

The key point is that safety, transfer, and scalability are not separate from
the cognitive-agent agenda; they determine whether that agenda can matter in real
systems.

% ==========================================================
\section{Cross-Cutting Synthesis}
\label{sec:crosscut}

\subsection{Cross-Paradigm Master Comparison}

Academic surveys usually benefit from a consolidated master table that lets
readers compare representative methods side by side rather than reconstructing
the landscape from dispersed narrative descriptions. Table~\ref{tab:master_methods}
therefore complements the family-level summary in Table~\ref{tab:families} with
a method-level comparison spanning the three paradigms. The goal is not to list
every paper reviewed in this survey, but to expose the most influential design
choices: what each method optimizes, what scenario it was validated on, whether
it relies on communication or explicit safety structure, and which limitation
remained unsolved after it appeared.

\begin{table}[ht]
\centering
\caption{Master comparison of representative methods across paradigms. The
I–II label marks bridge methods that emerged before full MARL became
dominant but already made interaction structure explicit.
\label{tab:master_methods}}
\setlength{\tabcolsep}{2pt}
\footnotesize
\begin{tabular*}{\textwidth}{@{\extracolsep{\fill}}p{1.6cm} p{0.7cm} p{1.15cm} p{1.75cm} p{1.6cm} p{1.45cm} p{1.45cm} p{1.7cm} p{2.0cm}}
\toprule
\textbf{Method} & \textbf{Year} & \textbf{Paradigm} & \textbf{Core mechanism} &
\textbf{Scenario / benchmark} & \textbf{Safety strategy} &
\textbf{Comm. dependency} & \textbf{Key strength} &
\textbf{Primary limitation} \\
\midrule
DAVE-2~\cite{bojarski2016end} & 2016 & I & End-to-end behavior cloning from raw pixels &
Real-road lane keeping & Implicit through expert data & None &
Established end-to-end driving feasibility &
Severe covariate shift and weak interaction reasoning \\

ChauffeurNet~\cite{bansal2018chauffeur} & 2018 & I & Perturbation-augmented imitation with auxiliary losses &
Urban driving logs & Collision/off-road penalties during training & None &
Learns recovery behavior beyond nominal demonstrations &
Still lacks explicit multi-agent game structure \\

Social LSTM~\cite{alahi2016social} & 2016 & I–II & Grid-based social pooling &
Pedestrian interaction forecasting & Social collision-avoidance prior & None &
First influential explicit interaction encoder &
Grid scaling and limited map structure \\

VectorNet~\cite{gao2020vectornet} & 2020 & II & Vectorized scene graph with hierarchical interaction modeling &
Argoverse forecasting & Better structural priors & None &
Efficient long-range map-agent reasoning &
Prediction and planning remain partly decoupled \\

PiP~\cite{song2020pip} & 2020 & II & Planning-informed prediction conditioned on ego candidates &
Interactive planning / nuScenes-style settings & Joint safety-efficiency selection &
None & Couples forecast and plan more realistically &
Still approximates full strategic interaction with candidate sets \\

QMIX~\cite{rashid2018qmix} & 2018 & II & Monotonic value decomposition under CTDE &
Cooperative MARL / SMAC-style tasks & Reward design rather than explicit hard constraints &
None at execution & Canonical cooperative MARL baseline &
Restricted expressiveness and mainly discrete-action focus \\

MADDPG~\cite{lowe2017multi} & 2017 & II & Centralized critic with decentralized actors &
Mixed cooperative-competitive MARL & Stabilized training under non-stationarity &
Train-time global state & Strong baseline for continuous multi-agent control &
Critic scaling degrades with many agents \\

MAPPO~\cite{yu2021surprising} & 2021 & II & PPO-style CTDE with centralized value function &
General cooperative benchmarks & Stable policy optimization & None at execution &
Simple, strong, reproducible baseline &
Limited explicit credit assignment machinery \\

V2VNet~\cite{wang2020v2vnet} & 2020 & II & Intermediate-feature V2V fusion with graph aggregation &
Cooperative perception in occluded scenes & Multi-view robustness through feature sharing &
V2V required & Realistic bridge from MARL communication to connected driving &
Sensitive to packet loss, synchronization, and bandwidth \\

MAT~\cite{wen2022multi} & 2022 & III & Autoregressive Transformer for joint multi-agent action modeling &
General MARL benchmarks & Implicit via sequence-level policy structure &
None at execution & Scales attention over agents and histories &
Safety and deployment realism are still external to the model \\

DiLu~\cite{wen2023dilu} & 2023 & III & LLM reasoning, reflection, and memory retrieval &
Closed-loop urban driving & Reflective error correction & Optional cloud / high compute &
Adapts to novel cases without retraining &
Reliability depends on LLM faithfulness and latency \\

LanguageMPC~\cite{sha2023languagempc} & 2023 & III & Slow-time language reasoning plus fast MPC control &
Closed-loop driving tasks & MPC hard constraints with semantic planning &
Optional cloud / high compute & Strong hybrid template for cognition plus control &
Formal guarantees for the full stack remain incomplete \\
\bottomrule
\end{tabular*}
\end{table}

\noindent Table~\ref{tab:master_methods} highlights a pattern that is easy to
miss when papers are discussed only within their local subfields. Landmark
methods do not differ merely in model size or benchmark score; they differ in
which \emph{problem formulation} they prioritize. DAVE-2 established the
feasibility of end-to-end driving, while ChauffeurNet directly targets
robustness to distributional shift. Social LSTM, VectorNet, and
PiP progressively reframe the core bottleneck as one of explicit interaction
representation and prediction-planning coupling. QMIX, MADDPG, and MAPPO then
treat coordination under non-stationarity as the central challenge. V2VNet
treats observability itself as a distributed systems problem, while MAT explores
how sequence modeling can scale coordination over larger agent contexts. DiLu
and LanguageMPC finally treat rare-event reasoning and semantic interpretation
as first-class bottlenecks. This is why a cross-paradigm table is especially
important in a survey of this field. With that method-level comparison in place,
we can turn from side-by-side description to the more selective question of
which methods actually became the enduring reference points for subsequent work.

\subsection{Which Methods Matter Most in Practice?}

Across the literature, a stable practical hierarchy has emerged. For
single-agent learning, the field-defining line is DAVE-2 $\rightarrow$
conditional imitation $\rightarrow$ ChauffeurNet
\cite{bojarski2016end,codevilla2018end,bansal2018chauffeur}. For generic
cooperative MARL, the most influential algorithmic anchors remain QMIX,
MADDPG, and MAPPO \cite{rashid2018qmix,lowe2017multi,yu2021surprising}. For
interaction-aware scene modeling, VectorNet and closely related graph or
polyline approaches became the main reference family
\cite{gao2020vectornet,liang2020lanegcn,salzmann2020trajectron}. For explicit
cooperative perception, V2VNet is more representative of realistic driving than
many generic communication papers \cite{wang2020v2vnet}. For the cognitive turn,
DiLu, GPT-Driver, and LanguageMPC currently define the design space more clearly
than dozens of loosely similar prompt-engineering variants
\cite{wen2023dilu,mao2023gpt,sha2023languagempc}.

This weighted view does not imply that the remaining papers are unimportant.
Rather, it indicates where the literature changed the field's shared vocabulary.
The most useful survey compression therefore preserves depth for these anchors
and moves adjacent variants into comparative clusters.

\subsection{Three Persistent Tensions}

\textbf{Credit assignment versus tractability.}
Value decomposition is attractive because it preserves decentralized action
selection, but its structural constraints can limit expressiveness. Actor-critic
methods are more flexible, especially in continuous spaces, but can suffer from
higher variance and weaker interpretability
\cite{sunehag2018value,rashid2018qmix,foerster2018counterfactual,lowe2017multi,peng2021facmac}.

\textbf{Explicit communication versus implicit coordination.}
CTDE often yields strong coordination without online communication, especially
when the shared reward already aligns agents. Explicit communication helps when
information is truly distributed, as in cooperative perception, but it adds
latency, bandwidth, and synchronization problems
\cite{das2019tarmac,sukhbaatar2016learning,wang2020v2vnet,wang2020bandwidth}.

\textbf{Closed-world optimization versus open-world robustness.}
Benchmark gains do not guarantee rare-event competence. This is the shared
failure that links the end of Paradigm~II to the rise of Paradigm~III. LLM-based
agents improve semantic flexibility, but only hybrids currently offer a plausible
path to both robustness and real-time control
\cite{caesar2021nuplan,fu2023drive,wen2023dilu,sha2023languagempc}.

\begin{table}[ht]
\centering
\caption{High-level comparison of the main method families reviewed in this survey.
\label{tab:families}}
\renewcommand{\arraystretch}{1.15}
\small
\begin{tabular}{p{2.7cm}p{2.0cm}p{2.0cm}p{2.2cm}p{5.0cm}}
\toprule
\textbf{Family} & \textbf{Action space} & \textbf{Coordination mechanism} & \textbf{Interpretability} & \textbf{Best fit in AD} \\
\midrule
Behavior cloning / imitation \cite{bojarski2016end,bansal2018chauffeur,codevilla2018end} & Continuous & None or implicit & Low--medium & Direct control under narrow distribution support \\
Value decomposition \cite{sunehag2018value,rashid2018qmix,son2019qtran} & Mostly discrete & Shared value factorization & Medium & Cooperative control with explicit team reward \\
Actor-critic CTDE \cite{lowe2017multi,foerster2018counterfactual,yu2021surprising,peng2021facmac} & Discrete or continuous & Centralized critic & Medium & Vehicle-level interaction and general MARL baselines \\
Communication-aware MARL \cite{sukhbaatar2016learning,das2019tarmac,wang2020v2vnet} & Both & Learned messages / feature sharing & Medium & Cooperative perception and information-limited coordination \\
Graph-based interaction models \cite{gao2020vectornet,liang2020lanegcn,salzmann2020trajectron} & N/A front-end & Relational scene encoding & Medium & Prediction, planning context, MARL front-ends \\
Transformer / LLM agents \cite{wen2022multi,wen2023dilu,mao2023gpt,sha2023languagempc} & Semantic or hybrid & Attention / high-level reasoning & High at the semantic level & Long-tail reasoning, explanation, hybrid tactical planning \\
\bottomrule
\end{tabular}
\end{table}

% ==========================================================
\section{Open Challenges and Future Directions}
\label{sec:challenges}

This section intentionally pairs each open challenge with the most plausible
near-term research direction, rather than listing unsolved problems in
isolation. That distinction matters in this literature: some bottlenecks are
primarily algorithmic (for example, credit assignment), some are infrastructural
(long-tail scenario generation and evaluation), and some are systems-level
(communication reliability and real-time LLM deployment). Treating them under a
single heading keeps the discussion concise while still preserving the survey
convention of moving from diagnosis to forward-looking synthesis.

\subsection{Challenge 1: Scaling Coordination Beyond Small Teams}

Most canonical MARL papers still evaluate on modest agent populations or highly
structured cooperation. Real traffic can involve dense, heterogeneous, partially
connected swarms. Scaling requires both algorithmic abstraction and benchmark
discipline. Mean-field and attention-style approximations are natural directions,
but they must be validated in driving-specific settings rather than generic MARL
alone \cite{zhou2020smarts,hernandez2019survey,li2020deep,wen2022multi}.

\subsection{Challenge 2: Credit Assignment Under Sparse Safety Signals}

Traffic rewards are sparse where it matters most: catastrophic failures are rare,
but decisive. Counterfactual baselines, Shapley-style decompositions, and
variance-aware constrained optimization remain promising, yet none has become a
universal solution for closed-loop AD \cite{foerster2018counterfactual,wang2020shapley,su2021vdac,prashanth2016variance}.

\subsection{Challenge 3: Safety Guarantees for Hybrid Agents}

As AD stacks combine learned interaction models, constrained policies, and LLM
reasoning, the safety problem becomes compositional. Safe RL and constraint
shields provide partial tools \cite{achiam2017constrained,alshiekh2018safe,shalev2016safe},
but future work needs interfaces that can verify or bound the behavior of
hybrid language-plus-control systems rather than only isolated components.

\subsection{Challenge 4: Realistic Communication Under Systems Constraints}

Cooperative driving papers still often assume more reliable communication than
production V2X can provide. Future work should evaluate policies under packet
loss, delay, adversarial load, and asymmetric sensing, not just idealized
message passing \cite{wang2020bandwidth,wang2021vbc,wang2020v2vnet}. This is a
systems problem as much as a learning problem.

\subsection{Challenge 5: Long-Tail Coverage and World Modeling}

The literature increasingly agrees that logged datasets and standard simulators do
not cover the true deployment tail. A major research direction is therefore the
construction of richer generative or scenario-centric world models that can
support adversarial stress testing, interaction-conditioned rollouts, and
counterfactual safety analysis \cite{tobin2017domain,finn2017model,caesar2021nuplan,fu2023drive}. Progress here may matter more than yet another small improvement on a fixed benchmark.

\subsection{Challenge 6: Real-Time Cognitive Reasoning}

LLM-based agents are attractive precisely where classical policies are brittle,
but inference cost makes them difficult to place directly in real-time control
loops. Efficient deployment likely requires smaller distilled models, retrieval
and memory mechanisms, or tiered architectures in which language models reason
only when uncertainty or novelty exceeds a threshold
\cite{wen2023dilu,mao2023gpt,sha2023languagempc}. The core research question is
not whether LLMs can drive in principle, but how much cognition is necessary and
where it should sit in the decision stack.

% ==========================================================
\section{Conclusion}
\label{sec:conclusion}

The development of learning-based autonomous driving can be read as a sequence of
corrections to earlier simplifications. Single-agent imitation learning showed
that perception-to-action policies were possible, but interaction and covariate
shift limited robustness. MARL then supplied the language of coordination,
credit assignment, and communication, but its benchmark gains did not fully solve
long-tail brittleness or interpretability. Transformer and LLM-based agents now
promise broader context aggregation and semantic reasoning, yet they introduce
new problems of latency, grounding, and verification.

The most credible synthesis is therefore hybrid. Strong AD systems are likely to
combine structured scene encoders, MARL-compatible policy learning, explicit but
resource-aware communication, constrained safety mechanisms, and high-level
reasoning modules reserved for ambiguous or rare scenarios. That agenda is more
modest than claims of end-to-end autonomy through any single paradigm, but it is
also more consistent with what the literature has actually demonstrated.

% ==========================================================
\bibliographystyle{unsrtnat}
\bibliography{references}

\end{document}
